\documentclass[12pt,a4paper]{article}

\usepackage[a4paper,margin=2.54cm,top=2.8cm,bottom=2.8cm]{geometry}
\usepackage[T1]{fontenc}
\usepackage[utf8]{inputenc}
\usepackage{newtxtext}
\usepackage{newtxmath}
\usepackage{microtype}
\usepackage{setspace}
\usepackage{booktabs}
\usepackage{longtable}
\usepackage{needspace}
\usepackage{array}
\usepackage{tabularx}
\usepackage{multirow}
\usepackage{graphicx}
\usepackage{caption}
\usepackage{subcaption}
\usepackage{float}
\usepackage[section]{placeins}
\usepackage{amsmath}
\usepackage{enumitem}
\usepackage{fancyhdr}
\usepackage{lastpage}
\usepackage{pdflscape}
\usepackage[round,authoryear]{natbib}
\usepackage[colorlinks=true,linkcolor=blue,citecolor=blue,urlcolor=blue]{hyperref}
\usepackage{nameref}

\newcommand{\papertitle}{[Full Paper Title]}
\newcommand{\paperhead}{[Short Header Title]}
\newcommand{\paperauthor}{[Author Name]}
\newcommand{\paperaffiliation}{[Institution, City, Country]}
\newcommand{\paperemail}{email@example.com}

\setstretch{1.15}
\setlength{\parindent}{0.5in}
\setlength{\parskip}{0pt}
\setlength{\headheight}{14.5pt}
\setlength{\textfloatsep}{10pt plus 2pt minus 2pt}
\setlength{\intextsep}{8pt plus 2pt minus 2pt}
\setlist[itemize]{leftmargin=1.6em, itemsep=2pt, topsep=4pt}
\setlist[enumerate]{leftmargin=1.8em, itemsep=2pt, topsep=4pt}

\pagestyle{fancy}
\fancyhf{}
\fancyhead[L]{\small\textit{\paperhead}}
\fancyhead[R]{\small\thepage}
\renewcommand{\headrulewidth}{0.4pt}
\renewcommand{\footrulewidth}{0pt}

\fancypagestyle{titlepagestyle}{%
  \fancyhf{}%
  \fancyhead[L]{\small\textit{\paperhead}}%
  \renewcommand{\headrulewidth}{0.4pt}%
  \renewcommand{\footrulewidth}{0pt}%
}

\DeclareCaptionLabelSeparator{apasep}{\\}
\captionsetup{
  font=small,
  labelfont=bf,
  textfont=it,
  labelsep=apasep,
  justification=raggedright,
  singlelinecheck=false,
  skip=6pt
}
\captionsetup[table]{position=above}
\captionsetup[figure]{position=above}
\captionsetup[longtable]{
  font=small,
  labelfont=bf,
  textfont=it,
  labelsep=apasep,
  justification=raggedright,
  singlelinecheck=false,
  width=\textwidth
}
\setlength{\LTcapwidth}{\textwidth}
\setlength{\LTleft}{0pt}
\setlength{\LTright}{0pt}

\newcommand{\fignote}[1]{%
  \par\vspace{3pt}\noindent%
  \begin{minipage}{0.96\linewidth}\footnotesize\raggedright\textit{Note.}~#1\end{minipage}%
}
\newcommand{\tabnote}[1]{%
  \par\vspace{3pt}\noindent%
  \begin{minipage}{0.96\linewidth}\footnotesize\raggedright\textit{Note.}~#1\end{minipage}%
}
\newcommand{\placeholder}[1]{\textit{#1}}

\begin{document}
\emergencystretch=2em

\begin{titlepage}
\thispagestyle{titlepagestyle}
\vspace*{1.4cm}
\begin{center}
{\fontsize{18}{22}\selectfont\bfseries \papertitle\par}
\vspace{1.8cm}
{\large \paperauthor\textsuperscript{1,*}\par}
\vspace{0.30cm}
{\normalsize \textsuperscript{1} \paperaffiliation\par}
\vspace{0.15cm}
{\normalsize \textsuperscript{*}Corresponding author: \href{mailto:\paperemail}{\paperemail}\par}
\end{center}
\vspace{1.0cm}
\begin{center}
\textbf{Abstract}
\end{center}
\begin{center}
\begin{minipage}{0.92\textwidth}
\placeholder{Write a compact abstract here. This template preserves the same title-page style and abstract block as the source project, but all study-specific claims, statistics, and wording have been removed. A practical starting point is: context, gap, objective, design, core result, and contribution.}
\end{minipage}
\end{center}
\end{titlepage}

\clearpage
\setcounter{page}{1}

\section{Introduction}

\subsection{Research Context and Motivation}

\placeholder{Explain the real-world problem, why it matters, and why the topic warrants study. Cite foundational background where relevant \citep{placeholder_background}.}

\subsection{Conceptual Background}

\placeholder{Define the main constructs, theoretical framing, or disciplinary lens. Keep this subsection at the level of argument-building rather than results preview.}

\subsection{State of the Literature and Gaps}

\placeholder{Summarize prior work, then isolate the unresolved gap your study addresses.}

\subsection{The Present Study and Research Questions}

\placeholder{State the purpose of the paper and list the confirmatory or guiding research questions.}

\begin{enumerate}[label=RQ\arabic*.,leftmargin=2.8em]
  \item \placeholder{[Primary research question]}
  \item \placeholder{[Secondary research question]}
  \item \placeholder{[Comparative or mechanism question]}
  \item \placeholder{[Model-performance or validation question]}
\end{enumerate}

\section{Methods}

\subsection{Framework or Corpus Construction}

\placeholder{Describe how the study frame, corpus, taxonomy, dataset, or evidence base was constructed. Reference protocol decisions and inclusion logic here.}

\subsection{Construct Inventory and Measurement}

\placeholder{Explain how core constructs, annotations, scales, or coded dimensions were defined and measured \citep{placeholder_method}.}

\subsection{Primary Outcome or Scoring Procedure}

\placeholder{Describe the central dependent measure, rating pipeline, extraction logic, or synthesis rule used in the study.}

\subsection{Statistical Analysis}

\subsubsection{Primary Model}

\placeholder{Document the main inferential model, outcomes, predictors, and thresholds.}

\subsubsection{Grouping or Clustering Strategy}

\placeholder{Describe any grouping, clustering, categorization, or stratification procedure.}

\subsubsection{Resampling or Confidence Intervals}

\placeholder{Specify bootstrap, cross-validation, confidence interval logic, or uncertainty estimation.}

\subsubsection{Permutation or Sensitivity Tests}

\placeholder{Describe label shuffling, robustness checks, negative controls, or alternative specifications.}

\subsubsection{Supplementary Analyses}

\placeholder{List the extra exploratory or robustness analyses that support the primary results.}

\subsection{Code and Data Availability}

\placeholder{State where code, data, preregistration, and supporting assets will be archived or shared.}

\section{Results}

\subsection{Structure of the Research Space}

\placeholder{Summarize the first high-level descriptive result: composition, counts, map, taxonomy, corpus, sample, or framework structure.}

\begin{figure}[tb]
  \raggedright
  \caption{[Main Figure 1 Placeholder]}
  \label{fig:main1}
  \includegraphics[width=0.92\textwidth]{../assets/figures/fig1_framework_overview.pdf}
  \fignote{Replace this note with the explanatory note for your first main figure. Keep the note style if you want the same visual layout.}
\end{figure}

\subsection{Landscape and Group Profiles}

\placeholder{Describe the main structure in the latent space, map, sample segmentation, or grouped profile landscape.}

\begin{figure}[tb]
  \raggedright
  \caption{[Main Figure 2 Placeholder]}
  \label{fig:main2}
  \includegraphics[width=\textwidth]{../assets/figures/fig2_landscape.pdf}
  \fignote{Use this slot for the second main figure. The width and caption positioning match the source paper layout.}
\end{figure}

\subsection{Distribution and RQ1}

\placeholder{Report the overall distributional pattern that addresses the first research question.}

\begin{figure}[tb]
  \raggedright
  \caption{[Main Figure 3 Placeholder]}
  \label{fig:main3}
  \includegraphics[width=\textwidth]{../assets/figures/fig3_primary_heatmap.pdf}
  \fignote{Use this position for a matrix, heatmap, comparison grid, or other central results display.}
\end{figure}

\subsection{Primary Drivers and RQ2 and RQ4}

\placeholder{Report the primary inferential findings, including direction, uncertainty, and practical interpretation.}

\begin{table}[tb]
  \raggedright
  \small
  \caption{[Primary Table Placeholder]}
  \label{tab:primary_model}
  \resizebox{\textwidth}{!}{\input{../assets/tables/table3_model_results.tex}}
  \tabnote{Replace this note with the interpretation of the model table.}
\end{table}

\subsection{Group-Level Profiles and RQ3}

\placeholder{Describe how groups, clusters, strata, or categories differ on the main outcome.}

\subsection{Variance Structure, Latent Factors, and Network}

\placeholder{Summarize decomposition, factor, network, or structure-discovery results if applicable.}

\begin{figure}[tb]
  \raggedright
  \caption{[Main Figure 4 Placeholder]}
  \label{fig:main4}
  \includegraphics[width=\textwidth]{../assets/figures/fig4_network_or_structure.pdf}
  \fignote{This slot mirrors the last main figure position in the source paper.}
\end{figure}

\subsection{Validation or Convergence}

\placeholder{Report validation, convergence, inter-method agreement, ablation, or benchmark comparison results.}

\section{Discussion}

\subsection{Main Theoretical Interpretation}

\placeholder{Interpret the headline finding in relation to the theoretical framing from the introduction.}

\subsection{Protective or Moderating Factor}

\placeholder{Explain the strongest attenuating, moderating, or boundary-condition result.}

\subsection{Counterintuitive Pattern or Negative Finding}

\placeholder{Discuss a result that did not behave as expected and explain its implications.}

\subsection{Group-Specific Profiles and Practical Implications}

\placeholder{Translate the differentiated results into intervention, design, policy, or applied implications.}

\subsection{Scientific and Practical Value of the Approach}

\placeholder{Explain the value of the dataset, framework, scoring design, or broader methodological contribution.}

\subsection{Limitations}

\placeholder{List the most important validity constraints, assumptions, and external limitations.}

\subsection{Future Directions}

\placeholder{Explain the most defensible next-step studies or extensions.}

\section{Conclusion}

\placeholder{End with a tight summary of the contribution, core result, and why it matters.}

\clearpage
\bibliographystyle{apalike}
\bibliography{other/references}

\clearpage
\setcounter{figure}{0}
\renewcommand{\thefigure}{S\arabic{figure}}
\setcounter{table}{0}
\renewcommand{\thetable}{S\arabic{table}}

\section*{Supplementary Materials}

\subsection*{Supplementary Figures}

\begin{figure}[!htb]
  \raggedright
  \caption{[Supplementary Figure S1 Placeholder]}
  \label{fig:s1}
  \includegraphics[width=0.95\textwidth]{../assets/figures/figS1_workflow.pdf}
  \fignote{Use this slot for workflow, screening, review logic, or study process figures.}
\end{figure}

\begin{figure}[!htb]
  \raggedright
  \caption{[Supplementary Figure S2 Placeholder]}
  \label{fig:s2}
  \includegraphics[width=\textwidth]{../assets/figures/figS2_full_matrix.pdf}
  \fignote{Use this slot for a full matrix, expanded view, or diagnostic overview.}
\end{figure}

\begin{figure}[p]
  \raggedright
  \caption{[Supplementary Figure S3 Placeholder]}
  \label{fig:s3}
  \includegraphics[width=\textwidth]{../assets/figures/figS3_component_map.pdf}
  \fignote{Use this slot for component maps, facet views, or dimensional diagnostics.}
\end{figure}

\begin{figure}[!t]
  \raggedright
  \caption{[Supplementary Figure S4 Placeholder]}
  \label{fig:s4}
  \includegraphics[width=\textwidth]{../assets/figures/figS4_pca_or_latent_space.pdf}
  \fignote{Use this slot for PCA, embeddings, reduced-space diagnostics, or latent geometry plots.}
\end{figure}

\begin{figure}[!t]
  \raggedright
  \caption{[Supplementary Figure S5 Placeholder]}
  \label{fig:s5}
  \includegraphics[width=\textwidth]{../assets/figures/figS5_cluster_diagnostics.pdf}
  \fignote{Use this slot for cluster validation or model-selection diagnostics.}
\end{figure}

\clearpage
\begin{figure}[!t]
  \raggedright
  \caption{[Supplementary Figure S6 Placeholder]}
  \label{fig:s6}
  \includegraphics[width=\textwidth]{../assets/figures/figS6_similarity_structure.pdf}
  \fignote{Use this slot for similarity, factor, or structure-recovery diagnostics.}
\end{figure}

\begin{figure}[!t]
  \raggedright
  \caption{[Supplementary Figure S7 Placeholder]}
  \label{fig:s7}
  \includegraphics[width=\textwidth]{../assets/figures/figS7_validation.pdf}
  \fignote{Use this slot for validation scatterplots, agreement analysis, or benchmark comparison.}
\end{figure}

\begin{figure}[!t]
  \raggedright
  \caption{[Supplementary Figure S8 Placeholder]}
  \label{fig:s8}
  \includegraphics[width=\textwidth]{../assets/figures/figS8_non_additivity_or_robustness.pdf}
  \fignote{Use this slot for robustness, non-additivity, or higher-order diagnostic results.}
\end{figure}

\clearpage
\begin{figure}[!t]
  \raggedright
  \caption{[Supplementary Figure S9 Placeholder]}
  \label{fig:s9}
  \includegraphics[width=\textwidth]{../assets/figures/figS9_combined_panels.pdf}
  \fignote{Use this slot for multi-panel summary diagnostics.}
\end{figure}

\begin{figure}[!t]
  \raggedright
  \caption{[Supplementary Figure S10 Placeholder]}
  \label{fig:s10}
  \includegraphics[width=\textwidth]{../assets/figures/figS10_profiles.pdf}
  \fignote{Use this slot for within-group profiles, distributions, or response-pattern summaries.}
\end{figure}

\subsection*{Supplementary Tables}

\begin{table}[H]
  \raggedright
  \small
  \caption{[Supplementary Table S1 Placeholder]}
  \label{tab:s1}
  \resizebox{\textwidth}{!}{\input{../assets/tables/table1_summary.tex}}
  \tabnote{Replace with the summary table note.}
\end{table}

\begin{table}[htb]
  \raggedright
  \small
  \caption{[Supplementary Table S2 Placeholder]}
  \label{tab:s2}
  \resizebox{\textwidth}{!}{\input{../assets/tables/table2_key_rows.tex}}
  \tabnote{Replace with the note for ranked or key-row results.}
\end{table}

\begin{table}[htb]
  \raggedright
  \small
  \caption{[Supplementary Table S3 Placeholder]}
  \label{tab:s3}
  \resizebox{\textwidth}{!}{\input{../assets/tables/tableS3_regression_detail.tex}}
  \tabnote{Replace with the note for detailed coefficients or per-group models.}
\end{table}

\begin{table}[htb]
  \raggedright
  \small
  \caption{[Supplementary Table S4 Placeholder]}
  \label{tab:s4}
  \resizebox{\textwidth}{!}{\input{../assets/tables/tableS4_group_matrix.tex}}
  \tabnote{Replace with the note for the group-level matrix or comparison table.}
\end{table}

\begin{table}[htb]
  \raggedright
  \small
  \caption{[Supplementary Table S5 Placeholder]}
  \label{tab:s5}
  \resizebox{\textwidth}{!}{\input{../assets/tables/tableS5_group_profiles.tex}}
  \tabnote{Replace with the note for profile summaries and representative examples.}
\end{table}

\clearpage
\begingroup
\small
\raggedright
\input{../assets/tables/tableS6_inventory.tex}
\endgroup
\tabnote{Use this slot for a long inventory-style supplementary table.}

\clearpage
\begingroup
\footnotesize
\raggedright
\input{../assets/tables/tableS7_measurement_lexicon.tex}
\endgroup
\tabnote{Use this slot for a long measurement, lexicon, coding, or feature-inventory table.}

\clearpage
\begin{table}[H]
  \raggedright
  \small
  \caption{[Supplementary Table S8 Placeholder]}
  \label{tab:s8}
  \resizebox{\textwidth}{!}{\input{../assets/tables/tableS8_weights_or_loadings.tex}}
  \tabnote{Replace with the note for weights, loadings, sensitivity values, or component summaries.}
\end{table}

\begingroup
\small
\raggedright
\input{../assets/tables/tableS9_full_leaf_inventory.tex}
\endgroup
\tabnote{Use this slot for a final long-table inventory or appendix listing.}

\end{document}
